\begin{abstract}
Persistent-memory agents are often evaluated as if changing what they store implied changing how they later behave. We show that this inference can fail even under a controlled state intervention. Using the released ReasoningBank writer, we hold a source trajectory fixed and flip only the terminal success/failure reflection branch. Reward-conditioned writing is robust: all 20 complete Shopping pairs produce different persistent memories, and on a disjoint eight-trajectory control the between-branch memory distance exceeds a stronger semantic-preserving same-mode wording perturbation by 0.105 (exact $p=0.0078$). Yet the effect attenuates sharply across the native reuse pipeline. A forced fixed-evidence memory injection can produce terminal sensitivity (mean absolute success-rate difference $0.15625$, permutation $p=0.00074$), showing that the paired states can matter when directly supplied. Under native Shopping retrieval, however, branch-specific memory is exposed on 125/172 held-out retrieval opportunities while the terminal success difference falls to $0.02083$ ($p=0.4289$; 34/36 cells are zero), and the first-action total-variation difference is $0.06944$ ($p=0.5801$; 0/36 modal actions differ). A second-domain Reddit replication again shows robust write divergence (4/4) but sparse, sign-heterogeneous native terminal transport (mean absolute difference $0.125$, $p=0.2253$; 6/8 cells are zero). These results identify a stage-resolved boundary: \emph{persistent-state divergence is not behavioral divergence}. Reward-conditioned memory is a strong upstream state intervention, but realized behavioral authority is filtered by retrieval exposure and policy uptake. We therefore argue that self-improving memory systems should be evaluated as a transport chain---write $\rightarrow$ exposure $\rightarrow$ uptake $\rightarrow$ outcome---rather than by memory distance or endpoint performance alone.
\end{abstract}
